\section{Analytical Results}
\label{sec:results}

This section states and proves the three main propositions of the paper.
All results are derived analytically from the linearized model; quantitative
magnitudes are discussed in Section \ref{sec:quantitative}.

\subsection{Proposition 1: Endogenous Phillips Curve Flattening}

In the competitive benchmark, optimal Calvo price-setting yields the standard
forward-looking New Keynesian Phillips Curve (NKPC):
\begin{equation}
  \pi_t = \beta\,\mathbb{E}_t\pi_{t+1} + \kappa^{comp}\,\tilde{y}_t
  \qquad \text{where} \qquad
  \kappa^{comp} = \kappa_{mc}\!\left(\frac{1}{\sigma} + \frac{\varphi}{\alpha_n}\right)
  \label{eq:NKPC_comp}
\end{equation}
and $\kappa_{mc} \equiv (1-\xi)(1-\beta\xi)/\xi$ is the Calvo coefficient on
real marginal cost, $\sigma$ is the coefficient of relative risk aversion,
$\varphi$ is the inverse Frisch elasticity, and $\alpha_n \in (0,1)$ is the
labour cost share.

\begin{theorem}[Monopsony Phillips Curve Flattening]
  \label{thm:pc_flattening}
  In the presence of structural monopsony with countercyclical elasticity
  $\varepsilon_t = \bar\varepsilon - \gamma(u_t - \bar u)$ and $\gamma > 0$,
  the aggregate NKPC takes the same form as \eqref{eq:NKPC_comp} but with a
  strictly smaller slope:
  \begin{equation}
    \kappa^{mono} = \kappa_{mc}\!\left(\frac{1}{\sigma} + \frac{\varphi}{\alpha_n}
      - \underbrace{\frac{\gamma}{(\bar\varepsilon+1)^2(1-\bar u)\,\alpha_n}}_{\delta_\mu > 0}\right)
    < \kappa^{comp}
    \label{eq:kappa_mono}
  \end{equation}
  Moreover, $\kappa^{mono}$ is \emph{state-dependent}: $\kappa^{mono}$ falls
  further when unemployment exceeds its steady-state value.
\end{theorem}

\begin{proof}
  Real marginal cost in the monopsonistic economy is:
  \begin{equation}
    mc_t = \frac{w_t}{\mu_t^w \cdot z_t \cdot f_n(n_t)}
    \label{eq:mc_def}
  \end{equation}
  where $w_t$ is the nominal wage, $\mu_t^w = \varepsilon_t/(\varepsilon_t+1)$
  is the wage markdown, $z_t$ is TFP, and $f_n$ is the marginal product of labour.
  Log-linearising around the non-stochastic steady state (where $mc = 1$ and
  hats denote log-deviations):
  \begin{equation}
    \hat{mc}_t = \hat{w}_t - \hat{\mu}_t^w - \hat{z}_t - \widehat{f_n}_t
    \label{eq:mc_loglin}
  \end{equation}

  \paragraph{Markdown deviation.}
  The markdown satisfies $\mu^w = \varepsilon/(\varepsilon+1)$, so
  $\partial\mu^w/\partial\varepsilon = 1/(\varepsilon+1)^2 > 0$.
  With $\partial\varepsilon/\partial u = -\gamma$ and using the approximation
  $\hat u_t \approx -(1-\bar u)^{-1}\hat n_t$:
  \begin{equation}
    \hat\mu_t^w = \frac{1}{(\bar\varepsilon+1)^2}\,\partial\varepsilon/\partial u\cdot
      \hat u_t = \frac{\gamma}{(\bar\varepsilon+1)^2(1-\bar u)}\,\hat n_t
    \label{eq:mu_loglin}
  \end{equation}
  When employment rises ($\hat n_t > 0$), $\hat\mu_t^w > 0$ (workers gain
  bargaining power) and the markdown falls.

  \paragraph{Marginal product.}
  From Cobb-Douglas $f(n) = n^{\alpha_n}$:
  $\widehat{f_n}_t = (\alpha_n - 1)\hat n_t$.
  Aggregate output satisfies $\hat y_t = \alpha_n \hat n_t + \hat z_t$, so
  $\hat n_t = (\hat y_t - \hat z_t)/\alpha_n$ and
  $\widehat{f_n}_t = -[(1-\alpha_n)/\alpha_n](\hat y_t - \hat z_t)$.

  \paragraph{Wage.}
  The household labour supply FOC gives the competitive wage:
  $\hat w_t^{comp} = \sigma^{-1}\hat c_t + \varphi\hat n_t$.
  With market clearing $\hat c_t = \hat y_t$ and $\hat n_t = \hat y_t/\alpha_n$
  (at $\hat z_t = 0$ for exposition):
  $\hat w_t = (\sigma^{-1} + \varphi/\alpha_n)\hat y_t$.

  \paragraph{Combining.}
  Substituting into \eqref{eq:mc_loglin} and collecting terms in $\hat y_t$:
  \begin{align}
    \hat{mc}_t &= \underbrace{\left(\frac{1}{\sigma} + \frac{\varphi}{\alpha_n}\right)}_{\text{standard}}
      \hat y_t
      - \underbrace{\frac{\gamma}{(\bar\varepsilon+1)^2(1-\bar u)\,\alpha_n}}_{\delta_\mu}\hat y_t
      \notag\\[4pt]
    &= \left(\frac{1}{\sigma} + \frac{\varphi}{\alpha_n} - \delta_\mu\right)\hat y_t
    \label{eq:mc_adjusted}
  \end{align}
  Multiplying by $\kappa_{mc}$ and substituting into the standard Calvo pricing
  condition $\pi_t = \beta\mathbb{E}_t\pi_{t+1} + \kappa_{mc}\hat{mc}_t$
  yields \eqref{eq:kappa_mono}.

  Strict inequality holds because $\delta_\mu > 0$ for all $\gamma > 0$,
  $\bar\varepsilon > 0$, $\bar u \in (0,1)$, $\alpha_n \in (0,1)$.

  State-dependence: $\kappa^{mono}$ falls as $\bar u$ rises (more recession),
  since $\partial\delta_\mu/\partial\bar u > 0$ reduces the effective
  slope further. This is because higher unemployment compresses markdowns more
  aggressively, increasing the dampening of MC responses.
\end{proof}

\begin{corollary}[Sacrifice Ratio]
  \label{cor:sacrifice_ratio}
  The sacrifice ratio---output cost per percentage point of
  disinflation---is:
  \begin{equation}
    SR^{mono} = \frac{1}{\kappa^{mono}} > \frac{1}{\kappa^{comp}} = SR^{comp}
    \label{eq:sacrifice_ratio}
  \end{equation}
  Central banks must induce larger recessions to achieve the same disinflation
  in monopsonistic economies.
\end{corollary}

\paragraph{Magnitudes.}
In our baseline calibration ($\bar\varepsilon = 4$, $\gamma = 0.8$,
$\alpha_n = 0.65$, $\sigma = 2$, $\varphi = 2$, $\xi = 0.75$, $\beta = 0.985$):
\begin{equation*}
  \delta_\mu = \frac{0.8}{25 \times 0.95 \times 0.65} \approx 0.052
\end{equation*}
This gives $\kappa^{comp} = 0.311$, $\kappa^{mono} = 0.307$---a
1.4\% flattening of the Phillips curve slope. The associated sacrifice ratios
are $SR^{comp} = 3.21$ and $SR^{mono} = 3.26$. These magnitudes, computed from
the fully solved HANK model, are modest with standard calibration. Two remarks
are in order. First, $\delta_\mu$ scales with $\gamma$; higher countercyclical
sensitivity---plausible in industries with low labour mobility---produces
meaningfully larger flattening. Second, and more importantly, the \emph{HANK
redistribution channel} generates additional amplification beyond the slope
estimate: by compressing the incomes of high-MPC workers in recessions, monopsony
reduces aggregate demand by more than a representative-agent model would predict,
enlarging the effective sacrifice ratio above the pure $1/\kappa^{mono}$ measure.
The household Jacobian $\mathbf{G}^{C,r}$ computed via the sequence-space fake
news algorithm shows that aggregate consumption responds $-0.137$ per unit
interest rate change in the HANK model versus $-0.495$ in the representative-agent
benchmark, reflecting how MPC heterogeneity redistributes the burden of monetary
tightening toward constrained workers.

\subsection{Proposition 2: Fragmentation Amplifies Supply Shocks}

\begin{theorem}[Network Amplification Under Fragmentation]
  \label{thm:fragmentation}
  Let $\hat{z}_{sr,t}$ be a TFP shock to sector-region $(s,r)$ at date $t$.
  (i) The aggregate GDP impact of this shock is:
  \begin{equation}
    \frac{d\log Y_t}{d\hat z_{sr,t}} = \lambda_{sr}(\phi_t)
    \label{eq:domar_impact}
  \end{equation}
  where $\lambda_{sr}(\phi_t)$ is the time-varying Domar weight.
  (ii) $\lambda_{sr}(\phi_t)$ is weakly increasing in $\phi_t$ for upstream
  sectors $(s,r)$ with significant cross-bloc linkages:
  $\partial\lambda_{sr}/\partial\phi > 0$.
  (iii) The aggregate TFP semi-elasticity with respect to the fragmentation
  premium is:
  \begin{equation}
    \frac{d\log\mathrm{TFP}}{d\phi} = -\!\!\!
    \sum_{\substack{(s,r),(s',r')\\b(r)\neq b(r')}}
    \lambda_{sr}\,\Omega_{(s,r)(s',r')}\,\frac{1}{\tau^{rr'}} < 0
    \label{eq:tfp_loss}
  \end{equation}
  (iv) Supply shocks have strictly longer persistence in fragmented networks
  because domestic substitution operates at the lower within-bloc elasticity
  $\theta < \eta$.
\end{theorem}

\begin{proof}
  Part (i): \citet{hulten1978} theorem extended to the multi-sector CES
  environment of \citet{baqaee2019}. The Hulten result states that, to a first
  order, the GDP impact of a sectoral TFP shock equals that sector's Domar
  weight $\lambda_{sr} = P_{sr}X_{sr}/(P\cdot\mathrm{GDP})$.

  Part (ii): Differentiate the Leontief inverse with respect to $\phi$ using
  the matrix identity
  $d\mathbf{L}/d\phi = \mathbf{L}\,(d\boldsymbol\Omega^\top/d\phi)\,\mathbf{L}$.
  From the Armington CES structure, a rise in $\phi$ raises the iceberg cost
  on cross-bloc links, so
  $\partial\Omega_{(s,r)(s',r')}/\partial\phi < 0$ for cross-bloc pairs. This
  redirects final demand toward domestic suppliers, increasing their row-sum
  in $\mathbf{L}$ and thus their Domar weight. For upstream sectors that
  supply intermediate goods across blocs, this reallocation is large; for
  downstream service sectors, it is negligible. See Appendix~\ref{app:proofs}
  for the full matrix derivation.

  Part (iii): Apply the envelope theorem to the CES cost function. For a
  marginal increase $d\phi$ in the cross-bloc trade cost:
  \begin{equation*}
    dm^{rr'}_{ss'} = -\frac{m^{rr'}_{ss'}}{\tau^{rr'}}\,d\phi
    \quad \Rightarrow \quad
    d\log Y = -\!\!\!\sum_{\text{cross-bloc}}\frac{\lambda_{sr}\,\Omega_{(s,r)(s',r')}}
    {\tau^{rr'}}\,d\phi
  \end{equation*}
  which gives \eqref{eq:tfp_loss} after dividing by $d\phi$. Strict negativity
  holds because all summands are non-negative.

  Part (iv): Under global integration, a cost shock is absorbed by substituting
  toward cross-bloc suppliers at the Armington elasticity $\eta$. Under
  fragmentation, this margin is blocked, confining adjustment to domestic
  suppliers at the lower within-sector elasticity $\theta < \eta$. Shock
  persistence therefore scales with $\eta/\theta > 1$. See
  Appendix~\ref{app:proofs}.
\end{proof}

\paragraph{Magnitudes.}
Using the BEA 2019 Benchmark Input-Output Table aggregated to 10 sectors and
5 regions (Appendix~\ref{app:data}), the calibrated formula \eqref{eq:tfp_loss}
gives $d\log\mathrm{TFP}/d\phi = -0.032$. A 15pp increase in $\phi$ thus
reduces aggregate TFP by 0.47\%. Upstream Domar weights (Energy, Mining,
Basic Manufacturing) increase by 15--25\% relative to baseline following such
a shock, consistent with the Domar reallocation mechanism of Part (ii).

\paragraph{Stagflationary Impulse.}
Proposition 2 implies that fragmentation shocks are \emph{stagflationary}: they
reduce TFP (deflationary for output) while simultaneously raising intermediate
input costs (inflationary for prices). Standard Taylor rules respond to this
by raising rates, which reduces output further without effectively curbing
the supply-side inflation, worsening the welfare outcome.

\subsection{Proposition 3: Optimal Augmented Taylor Rule}

\begin{theorem}[Optimal Monetary Policy Under Monopsony and Fragmentation]
  \label{thm:optimal_rule}
  Let the welfare criterion be the second-order approximation to household
  utility:
  \begin{equation}
    \mathcal{L} = \mathbb{E}\sum_{t=0}^\infty\beta^t\left[
    \mathrm{Var}[\pi_t] + \frac{\kappa^{mono}}{\varepsilon_p}\mathrm{Var}[\tilde y_t]
    + \frac{1}{\varepsilon_w\,\bar\mu^w}\mathrm{Var}[\hat\mu_t^w]
    + \frac{\alpha_m}{\eta\,\bar\tau}\mathrm{Var}[\hat\phi_t]\right]
    \label{eq:welfare_criterion}
  \end{equation}
  The welfare-optimal Taylor rule is:
  \begin{equation}
    i_t = r^* + \pi^* + \phi_\pi^*(\pi_t-\pi^*) + \phi_y^*\,\tilde y_t
          + \phi_\mu^*\,\hat\mu_t^w + \phi_\phi^*\,\hat\phi_t
    \label{eq:optimal_rule}
  \end{equation}
  with $\phi_\pi^* > \phi_\pi^{std}$ and $\phi_\phi^* < 0$.
\end{theorem}

\begin{proof}
  The welfare loss function \eqref{eq:welfare_criterion} has two additional
  terms relative to the standard \citet{woodford2003} loss:
  \begin{enumerate}[(i)]
    \item The \emph{labour wedge term}
      $(1/\varepsilon_w\bar\mu^w)\mathrm{Var}[\hat\mu_t^w]$: welfare is directly
      lowered by fluctuations in the wage markdown because households value
      the consumption-leisure margin and are not compensated for the
      productivity-wage gap.
    \item The \emph{fragmentation term}
      $(\alpha_m/\eta\bar\tau)\mathrm{Var}[\hat\phi_t]$: trade cost fluctuations
      impose deadweight losses proportional to the intermediate input share
      and inversely proportional to the Armington elasticity.
  \end{enumerate}

  \paragraph{Ramsey first-order conditions.}
  The planner chooses $(\phi_\pi, \phi_y, \phi_\mu, \phi_\phi)$ to minimise
  $\mathcal{L}$ subject to the reduced-form law of motion from combining the
  NKPC \eqref{eq:NKPC_comp} and aggregate demand equation. Let
  $\rho = \rho_\phi$ denote the AR(1) persistence of the fragmentation shock
  and $\psi_{mc}$ the cost-push pass-through coefficient. The reduced-form
  output and inflation responses per unit fragmentation shock are:
  \begin{align}
    \Phi_y(\boldsymbol\phi)   &= \frac{-\sigma^{-1}\phi_\phi - \psi_z}
      {(1-\rho) + \sigma^{-1}\phi_y + \sigma^{-1}(\phi_\pi-\rho)\kappa^{mono}\alpha_n^{-1}/(1-\beta\rho)}
    \label{eq:Phi_y}\\[4pt]
    \Phi_\pi(\boldsymbol\phi) &= \frac{\kappa^{mono}}{\alpha_n(1-\beta\rho)}\Phi_y
      + \frac{\kappa^{mono}\psi_{mc}}{1-\beta\rho}
    \label{eq:Phi_pi}
  \end{align}
  where $\psi_z = |d\log\mathrm{TFP}/d\phi|$ is the TFP sensitivity from
  Proposition~2. The welfare loss is $\mathcal{L} = \Phi_\pi^2 V_\phi + \lambda_y
  \Phi_y^2 V_\phi + \cdots$ where $V_\phi = \sigma_\phi^2/(1-\rho^2)$.

  \paragraph{Sign of $\phi_\phi^*$.}
  From \eqref{eq:Phi_y}--\eqref{eq:Phi_pi}: $\partial\Phi_y/\partial\phi_\phi > 0$
  (less negative $\phi_\phi$ deepens the output trough) and
  $\partial\Phi_\pi/\partial\phi_\phi < 0$ (more accommodation reduces inflation).
  The FOC $\partial\mathcal{L}/\partial\phi_\phi = 0$ sets the inflation-output
  trade-off at the social optimum. Since a fragmentation shock is a supply shock
  (both $\Phi_y < 0$ and $\Phi_\pi > 0$, unlike a demand shock where both
  move together), it generates an inflation-output tradeoff under the standard
  rule. The planner resolves this by choosing $\phi_\phi^* < 0$, partially
  accommodating the cost-push component to avoid excessive output sacrifice.
  The optimal magnitude $|\phi_\phi^*|$ is increasing in the welfare weight
  $\lambda_\phi = \alpha_m/(\eta\bar\tau)$ on fragmentation variance and in
  the TFP sensitivity $\psi_z$.

  \paragraph{Sign of $\phi_\pi^*$.}
  Because monopsony flattens the NKPC ($\kappa^{mono} < \kappa^{comp}$), the
  planner must use a larger $\phi_\pi^*$ to achieve the same inflation
  stabilisation.  The optimal augmented rule tilts the inflation-output frontier
  inward relative to the standard rule; a larger $\phi_\pi^*$ is part of
  reaching the efficient point. See Appendix~\ref{app:proofs} for the
  complete Ramsey derivation.
\end{proof}

\begin{corollary}[Welfare Gain from Augmented Rule]
  The welfare gain from switching from the optimal standard Taylor rule to the
  optimal augmented rule \eqref{eq:optimal_rule} is strictly positive and
  increasing in $\sigma_\phi^2$ (fragmentation shock variance).
\end{corollary}

\paragraph{Magnitudes.}
In our baseline calibration, the optimal augmented rule has
$\phi_\pi^* = 3.16$, $\phi_y^* = 0.00$, $\phi_\phi^* = -0.22$.
The negative coefficient on the fragmentation premium is the key prescription:
when $\hat\phi_t > 0$ (trade costs spike), the rule calls for \emph{less}
tightening than the standard Taylor rule would imply, absorbing the supply-side
cost-push rather than transmitting it into a demand contraction. Relative to
the optimal standard Taylor rule, the augmented rule reduces the total discounted
welfare loss by 0.06\% in baseline. This welfare gain scales with the variance of
fragmentation shocks: with $\sigma_\phi = 0.10$ (twice the baseline), the gain
rises to approximately 0.24\%, and with $\sigma_\phi = 0.20$, to roughly 1.0\%
of steady-state consumption. Given the historical volatility of supply chain
disruptions over 2018--2024, larger values of $\sigma_\phi$ are arguably
empirically relevant.
